\subsection{Аналіз результатів дослідно-експериментальної роботи}\label{sec:3.3}

\subsubsection{Результати констатувального етапу}

На констатувальному етапі (2020--2021) було проведено діагностику початкового рівня сформованості вмінь розробки ІОР серед студентів спеціальності 014 Середня освіта (Інформатика). Діагностика здійснювалася за критеріальним апаратом, описаним у підрозділі~\ref{sec:3.1}, з використанням всього комплексу діагностичних інструментів.

Результати вхідного анкетування засвідчили, що переважна більшість студентів мають лише побіжне уявлення про технології доповненої та віртуальної реальності. Зокрема, 82,8\% студентів вказали, що їхній досвід з AR обмежується використанням AR-фільтрів у соціальних мережах (Instagram, Snapchat); 14,1\% мали досвід використання AR-навігації або освітніх AR-додатків; лише 3,1\% мали будь-який досвід програмування, пов'язаний з AR/VR, і жоден студент не мав досвіду розробки WebAR додатків. При цьому 68,8\% студентів виявили зацікавленість у вивченні технологій AR/VR для освіти, що свідчить про наявність потенціалу для формування мотивації.

Розподіл студентів за рівнями сформованості вмінь на констатувальному етапі подано у табл.~\ref{tab:konst}.

\begin{table}[!h]
\caption{Розподіл студентів за рівнями сформованості вмінь розробки ІОР (констатувальний етап)}
\label{tab:konst}
\centering
\small
\begin{tabular}{|l|c|c|c|c|}
\hline
\hfill \textbf{Група} & \textbf{Низький, \%} & \textbf{Середній, \%} & \textbf{Достатній, \%} & \textbf{Високий, \%} \\
\hline
ЕГ ($n_1 = 34$) & 35,3 & 41,2 & 17,6 & 5,9 \\
\hline
КГ ($n_2 = 30$) & 33,3 & 43,3 & 16,7 & 6,7 \\
\hline
\end{tabular}
\end{table}

Як видно з табл.~\ref{tab:konst}, розподіли студентів ЕГ та КГ за рівнями сформованості вмінь на констатувальному етапі є подібними: в обох групах найбільша частка студентів має середній рівень (41,2\% в ЕГ та 43,3\% в КГ), значна частка~-- низький рівень (35,3\% в ЕГ та 33,3\% в КГ), лише незначна кількість~-- достатній (17,6\% та 16,7\%) і високий (5,9\% та 6,7\%) рівні.

Розподіл студентів ЕГ та КГ за рівнями на констатувальному етапі за кожним критерієм подано у табл.~\ref{tab:konst_detail}.

\begin{table}[!h]
\caption{Розподіл студентів за рівнями та критеріями (констатувальний етап, \%)}
\label{tab:konst_detail}
\centering
\small
\begin{tabular}{|l|l|c|c|c|c|}
\hline
\textbf{Критерій} & \textbf{Група} & \textbf{Низький} & \textbf{Середній} & \textbf{Достатній} & \textbf{Високий} \\
\hline
\multirow{2}{*}{Мотиваційний} & ЕГ & 29,4 & 41,2 & 23,5 & 5,9 \\
\cline{2-6}
 & КГ & 30,0 & 40,0 & 23,3 & 6,7 \\
\hline
\multirow{2}{*}{Когнітивний} & ЕГ & 35,3 & 47,1 & 11,8 & 5,9 \\
\cline{2-6}
 & КГ & 33,3 & 46,7 & 13,3 & 6,7 \\
\hline
\multirow{2}{*}{Операційно-діяльнісний} & ЕГ & 41,2 & 35,3 & 17,6 & 5,9 \\
\cline{2-6}
 & КГ & 40,0 & 36,7 & 16,7 & 6,7 \\
\hline
\multirow{2}{*}{Рефлексивно-оцінювальний} & ЕГ & 35,3 & 41,2 & 17,6 & 5,9 \\
\cline{2-6}
 & КГ & 33,3 & 43,3 & 16,7 & 6,7 \\
\hline
\end{tabular}
\end{table}

Аналіз табл.~\ref{tab:konst_detail} свідчить, що на початку експерименту найнижчі показники зафіксовано за операційно-діяльнісним критерієм (41,2\% студентів ЕГ мають низький рівень), що пояснюється відсутністю у студентів практичного досвіду розробки WebAR додатків. Відносно вищі показники за мотиваційним критерієм (23,5\% студентів ЕГ мають достатній рівень) свідчать про наявність зацікавленості в імерсивних технологіях.

Для перевірки статистичної еквівалентності ЕГ та КГ на констатувальному етапі було використано критерій $\chi^2$ Пірсона. Нульова гіпотеза $H_0$: розподіли студентів ЕГ та КГ за рівнями сформованості вмінь не мають статистично значущих відмінностей. Альтернативна гіпотеза $H_1$: розподіли мають статистично значущі відмінності.

Розрахунок емпіричного значення $\chi^2$ здійснювався за формулою:

\begin{equation}
\chi^2_{\text{емп}} = N \cdot M \cdot \sum_{i=1}^{c} \left( \frac{n_{1i}}{n_1} - \frac{n_{2i}}{n_2} \right)^2 \cdot \frac{1}{n_{1i} + n_{2i}}
\label{eq:chi2}
\end{equation}

\noindent де $N = n_1 + n_2 = 64$~-- загальна кількість студентів; $M = 2$~-- кількість груп; $c = 4$~-- кількість рівнів; $n_{1i}$, $n_{2i}$~-- кількість студентів $i$-го рівня в ЕГ та КГ відповідно; $n_1 = 34$, $n_2 = 30$~-- обсяги груп.

Для загального розподілу (табл.~\ref{tab:konst}) отримано $\chi^2_{\text{емп}} = 0{,}128$. Критичне значення $\chi^2_{\text{крит}} = 7{,}815$ при рівні значущості $\alpha = 0{,}05$ та кількості ступенів свободи $k = c - 1 = 3$. Оскільки $\chi^2_{\text{емп}} = 0{,}128 < \chi^2_{\text{крит}} = 7{,}815$, нульова гіпотеза $H_0$ не відхиляється, що підтверджує статистичну еквівалентність ЕГ та КГ на початку експерименту.

Аналогічну перевірку проведено для кожного критерію окремо; в усіх випадках $\chi^2_{\text{емп}} < \chi^2_{\text{крит}}$, що підтверджує еквівалентність груп за кожним критерієм.

%% =====================================================
%% §3.3.2 ФОРМУВАЛЬНИЙ ЕТАП (integrated quant + qual)
%% =====================================================

\subsubsection{Формувальний етап: динаміка навчального процесу}

Протягом формувального етапу (2021--2024) в експериментальній групі було впроваджено розроблену методику навчання розробки ІОР. Для аналізу перебігу формувального етапу використано два взаємодоповнюючі джерела: кількісний моніторинг (тестування, оцінювання лабораторних робіт, аналітика Moodle) та спрямований якісний контент-аналіз 14 транскриптів відеолекцій курсу <<Синтетична реальність в освіті>> (методику аналізу описано у підрозділі~\ref{subsec:3.2.qual}).

\subsubsubsection{Реалізація факторів CAMIL у навчальному процесі}

Аналіз педагогічного дискурсу за факторами когнітивно-афективної моделі іммерсивного навчання (CAMIL, підрозділ~\ref{subsec:1.1}) виявив виразну темпоральну динаміку їх реалізації у навчальному процесі. Розподіл кодованих сегментів за факторами CAMIL та етапами методики подано у табл.~\ref{tab:camil_stages}.

\begin{table}[!h]
\centering
\caption{Частота кодованих сегментів за факторами CAMIL та етапами методики}
\label{tab:camil_stages}
\small
\begin{tabular}{|l|c|c|c|c|}
\hline
\hfill \textbf{Фактор CAMIL} & \textbf{ТА} & \textbf{ПР} & \textbf{ІР} & \textbf{Усього} \\
 & (тижні 1--6) & (тижні 7--14) & (тижні 15--18) & \\
\hline
Мотивація & 26 & 21 & 15 & 62 \\
\hline
Тілесність & 27 & 66 & 15 & 108 \\
\hline
Когнітивне навантаження & 13 & 17 & 14 & 44 \\
\hline
Агентність & 19 & 10 & 15 & 44 \\
\hline
Саморегуляція & 15 & 16 & 7 & 38 \\
\hline
Присутність & 1 & 3 & 0 & 4 \\
\hline
\textbf{Усього} & \textbf{101} & \textbf{133} & \textbf{66} & \textbf{300} \\
\hline
\end{tabular}

\medskip
\footnotesize
\textit{Примітка.} ТА~-- теоретико-аналітичний етап, ПР~-- практико-розвивальний, ІР~-- інтегративно-рефлексивний. Один сегмент може бути кодований за кількома факторами.
\end{table}

Темпоральну динаміку факторів CAMIL протягом 18 тижнів курсу подано на рис.~\ref{fig:camil_temporal}.

\begin{figure}[!h]
\centering
\begin{tikzpicture}
\begin{axis}[
    width=\textwidth,
    height=8cm,
    xlabel={Тиждень курсу},
    ylabel={Частка сегментів (\%)},
    xmin=0.5, xmax=18.5,
    ymin=0, ymax=50,
    xtick={1,2,3,4,5,6,7,8,9,10,11,12,13,15,16,17,18},
    xticklabel style={font=\small},
    yticklabel style={font=\small},
    legend style={
        at={(0.5,-0.25)},
        anchor=north,
        legend columns=3,
        font=\small,
        draw=none,
    },
    grid=major,
    grid style={gray!20},
    % Stage boundary lines
    extra x ticks={6.5, 14.5},
    extra x tick labels={},
    extra x tick style={grid=major, grid style={black!50, dashed, line width=0.8pt}},
]

% Motivation (CAM-MOTI)
\addplot[thick, mark=triangle*, mark size=2pt, blue!80!black] coordinates {
    (1,10.0) (2,6.9) (3,6.3) (4,7.6) (5,7.6) (6,4.6)
    (7,4.6) (8,9.8) (9,7.4) (10,2.4) (11,6.9) (12,2.0)
    (13,9.3) (15,5.4) (16,5.3) (17,5.3) (18,4.5)
};
\addlegendentry{Мотивація}

% Cognitive load (CAM-COGN)
\addplot[thick, mark=square*, mark size=2pt, red!80!black] coordinates {
    (1,6.0) (2,1.7) (3,1.6) (4,3.8) (5,3.8) (6,3.7)
    (7,3.7) (8,7.8) (9,4.4) (10,7.3) (11,3.4) (12,8.0)
    (13,1.9) (15,7.1) (16,5.3) (17,5.3) (18,2.7)
};
\addlegendentry{Когнітивне навантаження}

% Embodiment (CAM-EMBO)
\addplot[thick, mark=*, mark size=2pt, green!60!black] coordinates {
    (1,8.0) (2,10.3) (3,19.0) (4,2.9) (5,2.9) (6,1.9)
    (7,1.9) (8,47.1) (9,39.7) (10,9.8) (11,5.2) (12,6.0)
    (13,9.3) (15,12.5) (16,2.3) (17,2.3) (18,4.5)
};
\addlegendentry{Тілесність}

% Agency (CAM-AGEN)
\addplot[thick, mark=diamond*, mark size=2pt, orange!80!black] coordinates {
    (1,2.0) (2,0.0) (3,4.8) (4,11.4) (5,11.4) (6,2.8)
    (7,2.8) (8,2.0) (9,8.8) (10,2.4) (11,0.0) (12,4.0)
    (13,0.0) (15,10.7) (16,4.5) (17,4.5) (18,2.7)
};
\addlegendentry{Агентність}

% Self-regulation (CAM-SREG)
\addplot[thick, mark=pentagon*, mark size=2pt, violet!80!black] coordinates {
    (1,4.0) (2,3.4) (3,6.3) (4,2.9) (5,2.9) (6,3.7)
    (7,3.7) (8,3.9) (9,10.3) (10,12.2) (11,0.0) (12,4.0)
    (13,0.0) (15,1.8) (16,3.0) (17,3.0) (18,1.8)
};
\addlegendentry{Саморегуляція}

% Presence (CAM-PRES)
\addplot[thick, mark=otimes*, mark size=2pt, cyan!60!black] coordinates {
    (1,2.0) (2,0.0) (3,0.0) (4,0.0) (5,0.0) (6,0.0)
    (7,0.0) (8,0.0) (9,0.0) (10,2.4) (11,3.4) (12,0.0)
    (13,0.0) (15,0.0) (16,0.0) (17,0.0) (18,0.0)
};
\addlegendentry{Присутність}

% Stage labels
\node[anchor=south, font=\footnotesize\itshape, text=black!60] at (axis cs:3.5,48) {Теоретико-аналітичний};
\node[anchor=south, font=\footnotesize\itshape, text=black!60] at (axis cs:10.5,48) {Практико-розвивальний};
\node[anchor=south, font=\footnotesize\itshape, text=black!60] at (axis cs:16.5,48) {Інтегративно-рефлексивний};

\end{axis}
\end{tikzpicture}
\caption{Динаміка факторів CAMIL у педагогічному дискурсі протягом 18 тижнів курсу}
\label{fig:camil_temporal}
\end{figure}


Аналіз табл.~\ref{tab:camil_stages} та рис.~\ref{fig:camil_temporal} дає змогу виявити характерні закономірності.

\textit{Мотивація} (62 сегменти) найбільш виражена на теоретико-аналітичному етапі (тижні~1--6), де викладач систематично демонструє можливості WebAR та формує ціннісне ставлення до імерсивних технологій. Зокрема, на першому занятті (тиждень~1) спостерігається найвища концентрація мотиваційних висловлювань (10\% сегментів):

\begin{quote}
\textit{<<...що ми можемо назвати дійсно доповненою реальністю. Давайте спробуємо зазирнути \textnormal{[...]} навіть якщо ви наразі не маєте досвіду, деякі речі є достатньо інтуїтивними. Давайте спробуємо зазирнути у посібник, адже інтуїція не повинна замінювати інформацію>>} (тиждень~1, 00:02:44).
\end{quote}

Це узгоджується з мотиваційним компонентом готовності (підрозділ~\ref{subsec:2.2}): створення ситуацій успіху та демонстрація реальних застосувань WebAR формують стійку внутрішню мотивацію.

\textit{Тілесність} (108 сегментів~-- найвищий показник) концентрується на практико-розвивальному етапі (тижні~7--14) з виразним піком на тижнях~8--9, де вивчається Face Tracking: 47,1\% сегментів тижня~8 та 39,7\% тижня~9 містять коди тілесності. Це зумовлено технологічною специфікою: робота з 468 опорними точками обличчя, накладання масок, перемикання камер безпосередньо залучають тілесний досвід студентів:

\begin{quote}
\textit{<<Канонічна модель обличчя \textnormal{[...]} яка має досить велику кількість точок. Будь ласка, знайомтеся~-- канонічна модель обличчя з чотирьохсот шістдесяти восьми точками>>} (тиждень~8, 00:13:56).
\end{quote}

\textit{Когнітивне навантаження} (44 сегменти) розподіляється відносно рівномірно за етапами, що свідчить про послідовне застосування стратегій управління когнітивним навантаженням протягом усього курсу. Це відповідає принципу поступового ускладнення, закладеному у методику (підрозділ~\ref{subsec:2.3}).

\textit{Агентність} (44 сегменти) найвища на теоретико-аналітичному етапі (19 сегментів), що може бути пов'язано з більшою кількістю моментів вибору на початку курсу (вибір середовища розробки, знайомство з альтернативами), та зростає знову на інтегративно-рефлексивному етапі (15 сегментів), коли студенти отримують більшу автономію у підсумкових проєктах.

\textit{Саморегуляція} (38 сегментів) присутня на всіх етапах з помітним піком на тижні~10 (12,2\% сегментів), що збігається з лабораторною роботою~2, яка вимагає складнішої рефлексії. Зниження на інтегративно-рефлексивному етапі (7 сегментів) є несподіваним та може пояснюватися тим, що саморегуляція на цьому етапі реалізується переважно через письмові форми (есе-рефлексії, взаємооцінювання), які не фіксуються у транскриптах лекцій.

\textit{Присутність} (4 сегменти) зафіксована лише епізодично, що є очікуваним: відчуття присутності виникає безпосередньо при тестуванні AR-додатків на мобільних пристроях, а не у лекційному дискурсі.

\subsubsubsection{Розвиток компонентів TPACK у педагогічному дискурсі}

Аналіз транскриптів за компонентами TPACK (підрозділ~\ref{subsec:1.1}) виявив виразне домінування технологічних знань (TK) з поступовим розвитком педагогічних та інтеграційних компонентів. Розподіл за етапами подано у табл.~\ref{tab:tpack_stages}.

\begin{table}[!h]
\centering
\caption{Частота кодованих сегментів за компонентами TPACK та етапами методики}
\label{tab:tpack_stages}
\small
\begin{tabular}{|l|c|c|c|c|}
\hline
\hfill \textbf{Компонент} & \textbf{ТА} & \textbf{ПР} & \textbf{ІР} & \textbf{Усього} \\
\hline
TK (технологічні знання) & 155 & 109 & 80 & 344 \\
\hline
PK (педагогічні знання) & 6 & 1 & 5 & 12 \\
\hline
TPK (технолого-педагогічні) & \multicolumn{3}{c|}{\textit{виявлено інтерпретативно}} & $\approx$15 \\
\hline
TPACK-інтегрований & \multicolumn{3}{c|}{\textit{виявлено інтерпретативно}} & $\approx$8 \\
\hline
\end{tabular}

\medskip
\footnotesize
\textit{Примітка.} TPK та TPACK-інтегрований не піддаються надійному автоматизованому кодуванню, оскільки потребують інтерпретації контексту; частоти встановлено при інтерпретативній верифікації.
\end{table}

Домінування TK (344 сегменти) є закономірним: курс <<Синтетична реальність в освіті>> орієнтований передусім на формування технологічних компетентностей з WebAR. Проте динаміка TK~-- від 155 сегментів на теоретико-аналітичному етапі до 80 на інтегративно-рефлексивному~-- свідчить про поступове зменшення частки суто технічного дискурсу на користь більш інтегрованих форм.

PK у <<чистому>> вигляді зафіксовано рідко (12 сегментів), переважно у контексті обговорення педагогічного обґрунтування ІОР. Проте інтерпретативний аналіз виявив $\approx$15 сегментів з TPK~-- випадки, коли викладач обговорює, \textit{як} застосувати технологію для досягнення \textit{навчальних} цілей. Найбільш показовим прикладом є обговорення віртуальної хімічної лабораторії (anchem-ar), де технологічні рішення (вибір маркерів, анімація <<зливання>> реактивів) безпосередньо зумовлені педагогічними вимогами (моделювання хімічних реакцій для учнів). TPACK-інтегрований компонент ($\approx$8 сегментів) з'являється переважно на інтегративно-рефлексивному етапі, коли студенти розробляють підсумкові проєкти з педагогічним обґрунтуванням.

Виявлений дефіцит TPK у лекційному дискурсі підтверджує теоретичне положення про <<TPK-розрив>> у підготовці вчителів інформатики (підрозділ~\ref{subsec:1.1}). Водночас наявність 67,6\% студентів ЕГ, здатних обґрунтувати педагогічну доцільність ІОР (за даними оцінювання підсумкових проєктів), свідчить, що формування TPK відбувається не лише через лекційний дискурс, а й через практичну діяльність з проєктування навчальних ресурсів.

\subsubsubsection{Реалізація педагогічних умов}

Якісний аналіз транскриптів дав змогу виявити конкретні механізми реалізації кожної з чотирьох педагогічних умов, визначених у підрозділі~\ref{subsec:2.3}. Розподіл кодованих сегментів подано у табл.~\ref{tab:conditions_evidence}.

\begin{table}[!h]
\centering
\caption{Реалізація педагогічних умов у навчальному процесі (за даними контент-аналізу)}
\label{tab:conditions_evidence}
\small
\begin{tabularx}{\linewidth}{|l|c|c|c|c|X|}
\hline
\hfill \textbf{Умова} & \textbf{ТА} & \textbf{ПР} & \textbf{ІР} & \textbf{$\Sigma$} & \hfil \textbf{Провідні прояви} \\
\hline
Психолого-педагогічна & 10 & 24 & 11 & 45 & Підтримка при помилках, нормалізація труднощів, заохочення \\
\hline
Методична & 8 & 4 & 3 & 15 & Посилання на посібник, структурованість матеріалу \\
\hline
Технологічна & 2 & 27 & 18 & 47 & Налаштування ngrok, тестування на мобільних, live coding \\
\hline
Організаційна & 26 & 12 & 16 & 54 & Домашні завдання, терміни здачі, Moodle, структура курсу \\
\hline
\end{tabularx}
\end{table}

\textit{Психолого-педагогічна умова} (45 сегментів) найбільш виражена на практико-розвивальному етапі (24 сегменти), де студенти стикаються з найбільшими технічними труднощами. Характерним є систематичне нормалізування труднощів та створення ситуацій успіху:

\begin{quote}
\textit{<<...для того, щоб не було такої ситуації, що ми нібито відстаємо~-- насправді жодного відставання немає. Якщо вам здається, що ви запізнюєтесь \textnormal{[...]} це нормально>>} (тиждень~3, 00:01:37).
\end{quote}

\textit{Технологічна умова} (47 сегментів) концентрується на практико-розвивальному (27) та інтегративно-рефлексивному (18) етапах, де відбувається інтенсивна робота з WebXR Device API, TensorFlow.js та комерційними платформами. Наявність лише 2 сегментів на теоретико-аналітичному етапі свідчить, що на початку курсу технологічне забезпечення розглядається побіжно, а основне занурення відбувається при безпосередній практичній роботі.

\textit{Організаційна умова} (54 сегменти~-- найвищий показник) домінує на теоретико-аналітичному етапі (26 сегментів), де формується структура курсу, пояснюються правила, терміни здачі та система оцінювання:

\begin{quote}
\textit{<<...каталог саме для синтетичної реальності в освіті \textnormal{[...]} домашнє завдання на тиждень \textnormal{[...]} 4 одне завдання раз і друге завдання>>} (тиждень~8, 00:01:35).
\end{quote}

\textit{Методична умова} (15 сегментів) зафіксована переважно на теоретико-аналітичному етапі, коли викладач найбільш активно посилається на навчальний посібник та його структуру. Відносно низька частота пояснюється тим, що методичне забезпечення (посібник, відеолекції) функціонує переважно як \textit{фоновий} ресурс: студенти звертаються до нього самостійно, а не під час лекційного дискурсу.

Виявлення реалізації всіх чотирьох педагогічних умов у транскриптах підтверджує, що розроблена методика навчання забезпечує їх системне впровадження у практику навчального процесу, а не лише декларує їх теоретично.

\subsubsubsection{Стратегії викладання та їх динаміка}

На основі індуктивного аналізу педагогічного дискурсу виділено шість стратегій викладання, що застосовуються протягом курсу. Їх розподіл за етапами методики подано на рис.~\ref{fig:teaching_strategies}.

\begin{figure}[!h]
\centering
\begin{tikzpicture}
\begin{axis}[
    ybar stacked,
    width=0.85\textwidth,
    height=8cm,
    bar width=35pt,
    ylabel={Частка (\%)},
    ymin=0, ymax=105,
    symbolic x coords={Теоретико-аналітичний, Практико-розвивальний, Інтегративно-рефлексивний},
    xtick=data,
    xticklabel style={font=\small, text width=3.5cm, align=center},
    yticklabel style={font=\small},
    legend style={
        at={(0.5,-0.3)},
        anchor=north,
        legend columns=3,
        font=\small,
        draw=none,
    },
    grid=major,
    grid style={gray!20},
    ymajorgrids=true,
    xmajorgrids=false,
]

% Modeling (bottom layer)
\addplot+[ybar, fill=blue!70, draw=blue!90] coordinates {
    (Теоретико-аналітичний, 46.9)
    (Практико-розвивальний, 43.4)
    (Інтегративно-рефлексивний, 38.6)
};
\addlegendentry{Моделювання}

% Scaffolding
\addplot+[ybar, fill=green!60, draw=green!80] coordinates {
    (Теоретико-аналітичний, 4.8)
    (Практико-розвивальний, 7.0)
    (Інтегративно-рефлексивний, 6.7)
};
\addlegendentry{Скафолдинг}

% Questioning
\addplot+[ybar, fill=orange!60, draw=orange!80] coordinates {
    (Теоретико-аналітичний, 14.8)
    (Практико-розвивальний, 11.7)
    (Інтегративно-рефлексивний, 15.7)
};
\addlegendentry{Запитування}

% Feedback
\addplot+[ybar, fill=violet!50, draw=violet!70] coordinates {
    (Теоретико-аналітичний, 16.2)
    (Практико-розвивальний, 15.2)
    (Інтегративно-рефлексивний, 13.4)
};
\addlegendentry{Зворотний зв'язок}

% Troubleshooting
\addplot+[ybar, fill=red!50, draw=red!70] coordinates {
    (Теоретико-аналітичний, 14.8)
    (Практико-розвивальний, 20.6)
    (Інтегративно-рефлексивний, 21.3)
};
\addlegendentry{Вирішення проблем}

% Real-world connection
\addplot+[ybar, fill=cyan!50, draw=cyan!70] coordinates {
    (Теоретико-аналітичний, 2.6)
    (Практико-розвивальний, 2.2)
    (Інтегративно-рефлексивний, 4.3)
};
\addlegendentry{Зв'язок з практикою}

\end{axis}
\end{tikzpicture}
\caption{Розподіл стратегій викладання за етапами методики навчання (\%)}
\label{fig:teaching_strategies}
\end{figure}


\textit{Моделювання/демонстрація} є домінуючою стратегією на всіх етапах ($\approx$43\% усіх стратегічних сегментів), з поступовим зниженням від 46,9\% на теоретико-аналітичному до 38,6\% на інтегративно-рефлексивному етапі. Це зниження свідчить про поступовий перехід від <<покажу~-- повторіть>> до більш автономних форм навчальної діяльності. Характерним прикладом моделювання є live coding~-- демонстрація написання коду в реальному часі з коментуванням кожного кроку:

\begin{quote}
\textit{<<...дозволю собі зробити якийсь файл \textnormal{[...]} ось тут відповідно \textnormal{[...]} це буде нашою основою. Тепер спробуємо виконати дії, які нам необхідні для встановлення необхідної функціональності>>} (тиждень~15, 00:39:30).
\end{quote}

\textit{Вирішення проблем (troubleshooting)} зростає від 14,8\% на теоретико-аналітичному до 21,3\% на інтегративно-рефлексивному етапі. Це відображає збільшення складності технологій та зростання кількості нетривіальних помилок, з якими стикаються студенти при роботі з WebXR та TensorFlow.js. Наявність розлогих епізодів діагностики помилок у live coding є характерною рисою педагогічного дискурсу курсу.

\textit{Запитування} зберігає відносно стабільну частку ($\approx$14\%), що свідчить про систематичне залучення студентів до діалогу.

\textit{Зворотний зв'язок} поступово знижується від 16,2\% до 13,4\%, що може бути пов'язано з переходом до більш самостійної роботи студентів та зменшенням потреби у зовнішньому підкріпленні.

\textit{Скафолдинг} ($\approx$6\%) та \textit{зв'язок з практикою} ($\approx$3\%) є менш частотними, проте відіграють важливу роль: скафолдинг забезпечує зв'язок між заняттями (<<пам'ятаєте, як ми це робили раніше>>), а зв'язок з практикою~-- мотивацію через реальні застосування.

Загальна тенденція~-- від інструктор-центрованих (моделювання, зворотний зв'язок) до студент-центрованих (вирішення проблем, зв'язок з практикою) стратегій~-- відповідає послідовності етапів методики (підрозділ~\ref{subsec:2.3}), де роль студента поступово зростає від спостерігача до самостійного розробника.

\subsubsubsection{Кількісний моніторинг формувального етапу}

Поточний моніторинг, що здійснювався протягом формувального етапу, дав змогу зафіксувати позитивні тенденції за кожним критерієм.

\textit{За мотиваційним критерієм} аналіз активності у Moodle засвідчив, що 88,2\% студентів ЕГ систематично виконували практичні завдання у визначені терміни, 73,5\% брали активну участь у форумах курсу, а 41,2\% виконували додаткові (необов'язкові) завдання. Найбільший мотиваційний ефект мали: можливість негайного тестування на мобільних пристроях (91,2\%), демонстрація віртуальної хімічної лабораторії (83,3\%), лабораторна робота з AR-візитівкою (79,4\%). Ці дані узгоджуються з високою частотою мотиваційних факторів CAMIL у транскриптах першого етапу (26 сегментів).

\textit{За когнітивним критерієм} результати тестування засвідчили поступове зростання: середній бал від 38,5\% на початку до 62,7\% (тиждень~6), 71,4\% (тиждень~10), 74,8\% (тиждень~14), 78,3\% (тиждень~17). Найбільші труднощі~-- при вивченні WebXR Device API (67,1\%) та інтеграції TensorFlow.js (65,8\%). Транскрипти підтверджують цю тенденцію: зростання частки troubleshooting-сегментів на тижнях~11--17 (з 14,8\% до 21,3\%) збігається з переходом до найскладніших технологій.

\textit{За операційно-діяльнісним критерієм} лабораторна робота~1 (тиждень~6): 26,5\% <<відмінно>>, 47,1\% <<добре>>, 23,5\% <<задовільно>>, 2,9\% <<незадовільно>>. Лабораторна робота~2 (тиждень~10): 20,6\% <<відмінно>>, 44,1\% <<добре>>, 29,4\% <<задовільно>>, 5,9\% <<незадовільно>>. Творчі завдання: 52,9\% студентів ЕГ продемонстрували здатність до самостійного вирішення нестандартних завдань. Пік тілесності у транскриптах (тижні~8--9, Face Tracking) збігається з інтенсивним практикумом.

\textit{За рефлексивно-оцінювальним критерієм} якість взаємооцінювання зросла від 17,6\% обґрунтованих зауважень до 64,7\%. У транскриптах саморегуляція концентрується на тижнях~9--10 (піки 10,3\% та 12,2\%), що збігається з лабораторною роботою~2 та перехідним етапом від Face Tracking до WebXR.

%% =====================================================
%% §3.3.3 РЕЗУЛЬТАТИ КОНТРОЛЬНОГО ЕТАПУ
%% =====================================================

\subsubsection{Результати контрольного етапу}

\subsubsubsection{Кількісний аналіз}

На контрольному етапі (2024--2025) проведено підсумкову діагностику рівня сформованості вмінь розробки ІОР у студентів ЕГ та КГ. Розподіл за рівнями подано у табл.~\ref{tab:kontrol}.

\begin{table}[!h]
\caption{Розподіл студентів за рівнями сформованості вмінь розробки ІОР (контрольний етап)}
\label{tab:kontrol}
\centering\small
\begin{tabular}{|l|c|c|c|c|}
\hline
\hfill \textbf{Група} & \textbf{Низький, \%} & \textbf{Середній, \%} & \textbf{Достатній, \%} & \textbf{Високий, \%} \\
\hline
ЕГ ($n_1 = 34$) & 5,9 & 23,5 & 47,1 & 23,5 \\
\hline
КГ ($n_2 = 30$) & 20,0 & 43,3 & 26,7 & 10,0 \\
\hline
\end{tabular}
\end{table}

Порівняння результатів констатувального та контрольного етапів засвідчує суттєві зміни у розподілі студентів ЕГ: частка з низьким рівнем зменшилася від 35,3\% до 5,9\% (у 6 разів); з достатнім~-- зросла від 17,6\% до 47,1\% (у 2,7 рази); з високим~-- від 5,9\% до 23,5\% (у 4 рази). Загалом, частка студентів ЕГ з достатнім та високим рівнями зросла від 23,5\% до 70,6\%. У КГ аналогічний показник зріс від 23,4\% до 36,7\%.

Розподіл за критеріями на контрольному етапі подано у табл.~\ref{tab:kontrol_detail}.

\begin{table}[!h]
\caption{Розподіл студентів за рівнями та критеріями (контрольний етап, \%)}
\label{tab:kontrol_detail}
\centering
\small
\begin{tabular}{|l|l|c|c|c|c|}
\hline
\textbf{Критерій} & \textbf{Група} & \textbf{Низький} & \textbf{Середній} & \textbf{Достатній} & \textbf{Високий} \\
\hline
\multirow{2}{*}{Мотиваційний} & ЕГ & 5,9 & 17,6 & 47,1 & 29,4 \\
\cline{2-6}
 & КГ & 16,7 & 40,0 & 30,0 & 13,3 \\
\hline
\multirow{2}{*}{Когнітивний} & ЕГ & 5,9 & 29,4 & 41,2 & 23,5 \\
\cline{2-6}
 & КГ & 20,0 & 43,3 & 26,7 & 10,0 \\
\hline
\multirow{2}{*}{Операційно-діяльнісний} & ЕГ & 5,9 & 23,5 & 47,1 & 23,5 \\
\cline{2-6}
 & КГ & 23,3 & 43,3 & 23,3 & 10,0 \\
\hline
\multirow{2}{*}{Рефлексивно-оцінювальний} & ЕГ & 5,9 & 23,5 & 52,9 & 17,6 \\
\cline{2-6}
 & КГ & 20,0 & 46,7 & 26,7 & 6,7 \\
\hline
\end{tabular}
\end{table}

Динаміку змін за критеріями подано у табл.~\ref{tab:dynamics}.

\begin{table}[!h]
\caption{Динаміка змін за критеріями сформованості вмінь (ЕГ, \%)}
\label{tab:dynamics}
\centering
\small
\begin{tabular}{|l|c|c|c|c|}
\hline
\hfill \textbf{Критерій (до / після)} & \textbf{Низький} & \textbf{Середній} & \textbf{Достатній} & \textbf{Високий} \\
\hline
Мотиваційний & 29,4 / 5,9 & 41,2 / 17,6 & 23,5 / 47,1 & 5,9 / 29,4 \\
\hline
Когнітивний & 35,3 / 5,9 & 47,1 / 29,4 & 11,8 / 41,2 & 5,9 / 23,5 \\
\hline
Операційно-діяльнісний & 41,2 / 5,9 & 35,3 / 23,5 & 17,6 / 47,1 & 5,9 / 23,5 \\
\hline
Рефлексивно-оцінювальний & 35,3 / 5,9 & 41,2 / 23,5 & 17,6 / 52,9 & 5,9 / 17,6 \\
\hline
\end{tabular}
\end{table}

Для перевірки статистичної значущості різниці між ЕГ та КГ використано критерій $\chi^2$ Пірсона. Нульова гіпотеза $H_0$: розподіли студентів ЕГ та КГ за рівнями не мають статистично значущих відмінностей. Результати подано у табл.~\ref{tab:chi2_summary}.

\begin{table}[!h]
\caption{Результати статистичної перевірки за критерієм $\chi^2$ Пірсона (контрольний етап)}
\label{tab:chi2_summary}
\centering
\small
\begin{tabular}{|l|c|c|c|l|}
\hline
\hfill \textbf{Розподіл} & $\chi^2_{\text{емп}}$ & $\chi^2_{\text{крит}}$ & $k$ &\hfill  \textbf{Рішення} \\
\hline
Загальний & 8,247 & 7,815 & 3 & $H_0$ відхилено \\
\hline
Мотиваційний & 8,563 & 7,815 & 3 & $H_0$ відхилено \\
\hline
Когнітивний & 8,102 & 7,815 & 3 & $H_0$ відхилено \\
\hline
Операційно-діяльнісний & 9,217 & 7,815 & 3 & $H_0$ відхилено \\
\hline
Рефлексивно-оцінювальний & 8,384 & 7,815 & 3 & $H_0$ відхилено \\
\hline
\end{tabular}

\medskip
\footnotesize
\textit{Примітка.} $\alpha = 0{,}05$; $k = c - 1 = 3$; $\chi^2_{\text{крит}}(0{,}05; 3) = 7{,}815$.
\end{table}

Емпіричні значення $\chi^2$ для загального розподілу та для кожного критерію перевищують $\chi^2_{\text{крит}} = 7{,}815$, що дає підстави відхилити $H_0$: відмінності між ЕГ та КГ є статистично значущими, розроблена методика є ефективнішою за традиційний підхід.

\subsubsubsection{Якісний аналіз за критеріями}

\textit{Вплив на технічні компетентності.} Аналіз підсумкових проєктів студентів ЕГ засвідчив: 76,5\% здатні самостійно розробити WebAR додаток з відстеженням зображень, 64,7\%~-- з відстеженням обличчя, 52,9\%~-- з відстеженням довкілля (WebXR), 44,1\%~-- з інтеграцією ML. У КГ: 40,0\%, 26,7\%, 13,3\%, 6,7\% відповідно. Градієнт від простіших до складніших типів відстеження є природним та відображає як градієнт складності технологій, так і послідовність їх вивчення (Three.js $\to$ MindAR Image $\to$ Face $\to$ WebXR $\to$ TensorFlow.js).

\textit{Вплив трикрокової методики ML-інтеграції.} Трикрокова методика~\cite{Semerikov2024WebARML_CoSinE, Semerikov2024WebAR_EduDim}: 73,5\% студентів ЕГ виконали крок~1 (інтеграція handpose), 58,8\%~-- крок~2 (Teachable Machine), 44,1\%~-- крок~3 (комплексна інтеграція). Транскрипти тижня~15 підтверджують інтенсивну роботу з TensorFlow.js та face-api.js, де частка troubleshooting-сегментів зростає до 21,3\%, що відображає об'єктивну складність ML-інтеграції.

\textit{Педагогічна компетентність.} 67,6\% студентів ЕГ здатні обґрунтувати педагогічну доцільність розробленого ІОР (проти 30,0\% у КГ). Це свідчить про формування не лише TK, а й елементів TPK та TPACK~-- попри те, що у лекційному дискурсі TPK зустрічається рідко, практичний досвід проєктування ІОР для конкретних предметних галузей формує інтеграційний компонент TPACK.

\textit{Типові труднощі}: налаштування ngrok (67,6\%), async/await (58,8\%), тривимірна графіка (52,9\%), WebXR Device API (47,1\%), TensorFlow.js (44,1\%). У транскриптах кожна з цих труднощів відображена у troubleshooting-епізодах, що свідчить про оперативне реагування на проблеми студентів.

\subsubsubsection{Тріангуляція кількісних та якісних результатів}

Відповідно до конвергентного паралельного дизайну~\cite{Creswell2018MixedMethods}, було проведено систематичне зіставлення кількісних та якісних даних для оцінки їх збіжності. Матрицю конвергенції подано у табл.~\ref{tab:convergence}.

\begin{table}[!h]
\centering
\caption{Матриця конвергенції кількісних та якісних результатів}
\label{tab:convergence}
\small
\begin{tabularx}{\linewidth}{|p{2.5cm}|p{3.5cm}|p{3.5cm}|X|}
\hline
\hfil \textbf{Критерій} & \textbf{\makecell{Кількісний\\ індикатор}} & \hfil \textbf{Якісне свідчення} & \hfil \textbf{Оцінка конвергенції} \\
\hline
Мотиваційний ($\chi^2=8{,}563$) &
Достатній+високий: 29,4\%$\to$76,5\% &
CAM-MOTI: 62 сегменти, найвищий на ТА етапі &
\textbf{Збіжність.} Мотиваційний дискурс на початку курсу передує кількісному зростанню \\
\hline
Когнітивний ($\chi^2=8{,}102$) &
Тест: 38,5\%$\to$78,3\%; достатній+високий: 17,7\%$\to$64,7\% &
CAM-COGN: 44 сегменти рівномірно; TK: 344 сегменти &
\textbf{Збіжність.} Систематичне управління когнітивним навантаженням відповідає зростанню тестових балів \\
\hline
Операційно-діяльнісний ($\chi^2=9{,}217$~-- найвищий) &
Достатній+високий: 23,5\%$\to$70,6\% &
TS-MODEL: 43\%, TS-SCAFF: 6\%; CAM-EMBO пік на тижнях~8--9 &
\textbf{Збіжність.} Домінування моделювання та тілесного залучення відповідає найвищій $\chi^2$ \\
\hline
Рефлексивно-оцінювальний ($\chi^2=8{,}384$) &
Взаємооцінювання: 17,6\%$\to$64,7\%; високий лише 17,6\% &
CAM-SREG: 38 сегментів, пік на тижні~10; мала присутність на ІР етапі &
\textbf{Часткова збіжність.} Саморегуляція у дискурсі нижча за очікувану; рефлексія формується через письмові форми \\
\hline
\end{tabularx}
\end{table}

Тріангуляція виявила такі ключові закономірності.

\textit{Конвергенція 1: Операційно-діяльнісний критерій.} Найвище значення $\chi^2_{\text{емп}} = 9{,}217$ (найбільша статистична значущість серед усіх критеріїв) конвергує з домінуванням стратегії моделювання/демонстрації ($\approx$43\% стратегічних сегментів) та скафолдингу ($\approx$6\%). Систематичний live coding з коментуванням кожного кроку~-- провідний механізм, який забезпечує формування практичних умінь. Пік тілесності на тижнях~8--9 (Face Tracking) додатково підтверджує, що залучення тілесного досвіду сприяє формуванню операційно-діяльнісної компетентності.

\textit{Конвергенція 2: Мотиваційний критерій.} Зростання достатнього та високого рівнів від 29,4\% до 76,5\% конвергує з високою частотою мотиваційних висловлювань на першому етапі (26 з 62 сегментів). Транскрипти демонструють, що мотиваційна робота здійснюється переважно на початку курсу, а далі підтримується через демонстрацію реальних результатів та зв'язок з практикою.

\textit{Конвергенція 3: Когнітивний критерій.} Зростання тестових балів (38,5\%$\to$78,3\%) конвергує з рівномірним розподілом когнітивного навантаження (13/17/14 сегментів за етапами). Це свідчить про те, що послідовне управління когнітивним навантаженням~-- не лише на початку, а протягом усього курсу~-- є ключовим фактором когнітивного зростання.

\textit{Часткова конвергенція: Рефлексивно-оцінювальний критерій.} Кількісне зростання (23,5\%$\to$70,5\%) є статистично значущим ($\chi^2 = 8{,}384$), проте рефлексивно-оцінювальний критерій має найнижчий показник високого рівня (17,6\% проти 23,5\% за іншими критеріями). Якісні дані пояснюють цей нюанс: саморегуляція у лекційному дискурсі зафіксована у 38 сегментах (найменше серед факторів CAMIL), причому на інтегративно-рефлексивному етапі~-- лише 7 сегментів. Це свідчить про те, що формування рефлексивних умінь відбувається переважно через письмові форми (есе-рефлексії, взаємооцінювання у Moodle), які не фіксуються у лекційних транскриптах. Таким чином, якісні дані \textit{доповнюють} кількісні, виявляючи механізм формування рефлексивної компетентності.

\textit{Виявлена розбіжність: TPACK та TPK.} Кількісні дані засвідчують, що 67,6\% студентів ЕГ здатні обґрунтувати педагогічну доцільність ІОР, проте у лекційних транскриптах TPK-сегменти зустрічаються рідко ($\approx$15 із $\approx$1000). Ця розбіжність є інформативною: формування TPK відбувається не стільки через лекційний дискурс, скільки через практичну діяльність з проєктування навчальних ресурсів, захист проєктів та есе-рефлексії. Таким чином, лекційний дискурс закладає TK-фундамент, а педагогічна інтеграція здійснюється через інші компоненти навчально-методичного забезпечення.

%% =====================================================
%% §3.3.4 РЕКОМЕНДАЦІЇ
%% =====================================================

\subsubsection{Рекомендації щодо вдосконалення освітньої практики}

На основі результатів дослідно-експериментальної роботи, включаючи кількісний аналіз результатів педагогічного експерименту та якісний контент-аналіз транскриптів лекцій, сформульовано рекомендації щодо вдосконалення навчання розробки імерсивних освітніх ресурсів майбутніх учителів інформатики.

\textbf{1. Рекомендації щодо добору технологій.} Для навчання WebAR рекомендовано використовувати: \textit{MindAR} як основну бібліотеку відстеження зображень та обличчя завдяки її стабільності, документованості, включенню TensorFlow.js та підтримці Three.js (на відміну від AR.js, що використовує простіші маркери $16 \times 16$ та не підтримує ML); \textit{Three.js} як основний фреймворк для тривимірного рендерингу (є основою для MindAR та більшості WebAR бібліотек); \textit{WebXR Device API} для відстеження довкілля як стандарт W3C; \textit{A-Frame} для швидкого прототипування та демонстрації концепцій; \textit{TensorFlow.js та Teachable Machine} для інтеграції машинного навчання. Комерційні платформи (8th Wall, Zappar) доцільно розглядати оглядово для ознайомлення з промисловими рішеннями.

\textbf{2. Рекомендації щодо проєктування курсу.} Послідовність опанування технологій має відповідати принципу зростання складності: основи Three.js та 3D-рендерингу $\to$ Image Tracking (MindAR) $\to$ Face Tracking (MindAR) $\to$ World Tracking (WebXR Device API) $\to$ ML Integration (TensorFlow.js). Кожна тема повинна супроводжуватися практичним завданням з негайним тестуванням на мобільному пристрої. Рекомендований обсяг~-- не менше 18 тижнів (один семестр). Кейс віртуальної хімічної лабораторії рекомендується як мотиваційний приклад та демонстрацію повного циклу розробки ІОР.

\textbf{3. Рекомендації щодо оцінювання.} Рекомендовано поєднувати: тестовий контроль теоретичних знань (40 питань за 10 розділами); оцінювання лабораторних робіт за визначеними критеріями; підсумковий проєкт з розробки повноцінного ІОР із педагогічним обґрунтуванням; процедуру взаємооцінювання для формування рефлексивних умінь; есе-рефлексії для аналізу власного досвіду.

\textbf{4. Рекомендації щодо інтеграції у Moodle.} Організація курсу у LMS Moodle у тижневому форматі з визначеними дедлайнами забезпечує систематичність навчання (88,2\% студентів виконували завдання вчасно). Рекомендовано використовувати: <<Завдання>> для здачі лабораторних робіт; <<Тест>> для контролю знань; <<Форум>> для взаємодопомоги; <<Семінар>> для взаємооцінювання.

\textbf{5. Рекомендації щодо інтеграції ML у навчання WebAR.} Трикрокова методика (готові ML-моделі $\to$ Teachable Machine $\to$ комплексна інтеграція) є ефективним підходом. Важливо підкреслювати освітній контекст ML-інтеграції: створення інтелектуальних імерсивних ресурсів, що адаптуються до дій користувача.

\textbf{6. Рекомендації щодо стратегій викладання.} Контент-аналіз транскриптів виявив, що домінуючою стратегією є моделювання/демонстрація ($\approx$43\%), зокрема live coding. Рекомендовано: забезпечити систематичне моделювання на початкових етапах з поступовим переходом до студент-центрованих стратегій; інтегрувати troubleshooting-епізоди як навчальний ресурс (помилки~-- не перешкоди, а навчальні ситуації); посилити рефлексивний компонент у лекційному дискурсі, оскільки аналіз виявив його дефіцит порівняно з іншими факторами CAMIL.
